\documentclass{rptuma}
% Shared preamble for the three structure proposals. Built against ../rptuma.cls by build.sh,
% so each proposal renders exactly as a thesis written in the department template would.
\usepackage{xcolor}
\usepackage{booktabs}
\usepackage{tabularx}
\usepackage{enumitem}

% The class loads its logo as img/RPTUlogo; the image lives one directory up.
\graphicspath{{../}}

\definecolor{materialgray}{gray}{0.38}
\newcommand{\material}[1]{\par{\raggedright\small\sffamily\color{materialgray}\textbf{Material:} #1\par}}
\newcommand{\repo}[1]{\path{#1}}
\newcommand{\us}{\,\textmu s}

\setlist{nosep,leftmargin=*}
% Repository paths and identifiers are long unbreakable words; loose spacing beats margin overruns.
\AtBeginDocument{\sloppy}
\setcounter{tocdepth}{2}
% KOMA appends a period to section numbers once appendices exist; the department sample has none.
\KOMAoptions{numbers=noendperiod}
\addto\captionsenglish{\renewcommand{\bibname}{Sources Used for This Proposal}}

\newdate{thesisDate}{07}{10}{2026}
\author{Author's Name}
\department{University of Kaiserslautern-Landau\\Department of Computer Science\\Distributed Systems Lab}
\supervisortitle{Examiner:}
\supervisor{Prof. Dr.-Ing. Jens B. Schmitt}
\secondsupervisor{M. Sc. Supervisor Name}

\newcommand{\frontmatternote}{%
\section*{Front matter (unchanged from the department template)}
\begin{itemize}
  \item \textbf{Title page} and \textbf{Eigenst\"andigkeitserkl\"arung}, as generated by \texttt{rptuma.cls}.
        The declaration requires disclosing generative-AI use beyond grammar and style.
  \item \textbf{Abstract} (English): the problem (deception at the web edge has an unknown cost), the
        approach (one mechanism on four edges, levelled to identical behaviour), the method (paired
        benchmark on three measurement planes) and the headline numbers.
  \item \textbf{Zusammenfassung}: the same abstract in German.
  \item \textbf{Contents}, \textbf{List of Figures}, \textbf{List of Tables}.
  \item \textbf{Glossary} (optional, the class's \texttt{glossar} environment): WADM, honeytoken, edge,
        tier, VU, iteration, SLO, cgroup, LuaJIT, PUC Lua, proxy-wasm.
\end{itemize}
In the department template's own example, websites are cited through an archived copy
(\texttt{web.archive.org}); do the same for any web source in the list at the end.}


\subject{Master Thesis -- Structure Proposal B}
\title{Template B: A Research-Question-Driven Empirical Study}

\begin{document}

\maketitle
\pagenumbering{roman}
\setcounter{page}{1}

\chapter*{About This Proposal}

\textbf{Organising idea.} The thesis is a measurement study. The system is described only as much
as the reader needs, and the body has one chapter per research question, each ending with an
explicit answer. This is the convention of empirical software engineering, with a study-design
chapter and categorised threats to validity \cite{wohlin2012experimentation}. Systems
characterisation papers use the same shape, with one section per factor studied
\cite{zhu2022dissecting}.

\textbf{Choose it if} the main contribution is the answers: what deception at the edge costs, and
whether it breaks anything. The three questions the benchmark was built around (performance,
latency, functionality; \texttt{BENCHMARK\_METHODOLOGY.md} \S6) map directly onto chapters, and a
fourth asks how much depends on the runtime.

\textbf{Trade-off.} Design and implementation shrink to one chapter (\emph{The System under
Study}). A reader looking for the engineering has to find it there or in the appendix.

\begin{center}
\small
\begin{tabularx}{\textwidth}{@{}rXr@{}}
\toprule
Ch. & Title & approx.\ pages \\
\midrule
1 & Introduction & 5 \\
2 & Background and Related Work & 12 \\
3 & The System under Study & 10 \\
4 & Study Design & 12 \\
5 & RQ1: Latency Added & 8 \\
6 & RQ2: CPU and Capacity & 12 \\
7 & RQ3: Functional Correctness & 8 \\
8 & RQ4: Influence of the Extension Runtime & 8 \\
9 & Discussion & 7 \\
10 & Conclusion & 3 \\
A--D & Appendices & as needed \\
\bottomrule
\end{tabularx}
\end{center}
{\footnotesize Page counts are a rough budget, not a department rule.}

\textbf{How to read the rest.} The table of contents is the proposed structure. Under each
heading is a short note on what to write there, and a \emph{Material} line pointing to the
place in the repository that already holds it (\repo{results/} and \repo{plots/} are short for
\repo{benchmarks/results/} and \repo{benchmarks/results/plots/}).

\frontmatternote

\tableofcontents
\cleardoublepage
\pagenumbering{arabic}

%======================================================================
\chapter{Introduction}
\section{Context}
Edge proxies can plant honeytokens into any web application without changing its code
\cite{kahlhofer2024application}. Whether that is affordable has not been measured.

\section{Research Gap}
Deception work studies whether a lure is effective; proxy-overhead work measures generic filters.
Nothing measures deception logic across edge runtimes.

\section{Research Questions}
\textbf{RQ1} How much latency does WADM add per request on each edge, and how does that change with
load? \textbf{RQ2} What does it cost in CPU and capacity, and where does that cost arise?
\textbf{RQ3} Does it break what the server did without it, and does it detect what it should?
\textbf{RQ4} How much of the answer depends on the extension runtime (LuaJIT, PUC Lua,
WebAssembly)?

\section{Contributions}
The study design, the four levelled implementations, the measured answers, and a public harness.

\section{Outline}
One line per chapter, naming which research question each answers.

%======================================================================
\chapter{Background and Related Work}
Background and related work combined, as is usual for a study of this kind; each section ends
with how this thesis differs.

\section{Honeytokens and Application-Layer Deception}
From honeypots to honeytokens and deceptive web artefacts, and how they are deployed without
changing code \cite{spitzner2003honeytokens,bowen2009baiting,juels2013honeywords,virvilis2014changing,kahlhofer2024application}

\section{Edge Proxies and Their Extension Runtimes}
Request phases, body buffering, and the runtimes compared: LuaJIT, PUC Lua 5.1, V8 WebAssembly.

\section{Measuring Proxy Overhead}
Service-mesh sidecar studies \cite{zhu2022dissecting}.

\section{Benchmarking Pitfalls}
Known flaws \cite{vanderkouwe2019sok,heiser_crimes} and coordinated omission; these motivate
Chapter~\ref{chap:B.design}.

\section{Research Gap Revisited}
Restate the gap against the work just reviewed, and show why each RQ addresses part of it.

%======================================================================
\chapter{The System under Study}
\label{chap:B.system}
Just enough of WADM to interpret the measurements.

\section{WADM in Brief}
Five honeytoken kinds plus the fake SQL-injection trap, all driven by one \texttt{config.json}. One
table: kind, where it is planted, what triggers it.
\material{\repo{CONTEXT.md}, \repo{config.json}}

\section{The Four Implementations}
A table: edge, runtime, hooks used, how bodies are buffered, what differs.
\material{\repo{nginx/README.md}, \repo{envoy_scripts/README.md}, \repo{wasm-filter/README.md},
\repo{apache_scripts/README.md}}

\section{Equivalence as a Precondition for Comparison}
The canonical mechanism, tuning parity (keep-alive, logging, code caching), and the parity check
that proves the four behave identically.
\material{\repo{docs/EDGE_LEVELING.md}, \repo{benchmarks/parity_check.py}}

%======================================================================
\chapter{Study Design}
\label{chap:B.design}
The experiment as a reproducible procedure.

\section{Variables}
Independent: edge, tier, arrival rate, request mix, page size. Dependent: latency difference, CPU
per request, capacity, share of intact responses, alert recall.

\section{Environment}
VM, cpuset pinning (system, edge, origin, load generator), guest tuning, image versions.
\material{\repo{docs/BENCHMARK_METHODOLOGY.md} \S4}

\section{Workload and Load Model}
An open-loop arrival rate of 1, 10 or 100 iterations/s; nine requests per iteration; probes in
Williams order \cite{williams1949designs}.
\material{\repo{benchmarks/test.js}}

\section{Instruments}
\subsection{In-Edge Timers}
Wall-clock regions around detection and injection, logged in \textmu s.
\subsection{End-to-End Latency}
The k6 summary, per probe.
\subsection{CPU per Request}
A host-side sampler of cgroup \texttt{cpu.stat}, integrated over the measured window.

\section{Paired Design and Replication}
Bare, silent and instrumented tiers in balanced order, 5 replicates. silent $-$ bare is the cost of
WADM; instrumented $-$ silent is the cost of measuring it.

\section{Validity Gates and Negative Control}
The \texttt{assess\_run()} thresholds; origin CPU as a negative control; detection of host freezes.
\material{\repo{benchmarks/wadm_timings.py}}

\section{Analysis Procedure}
Medians of paired differences, the range over replicates, ``resolved'' when the replicates agree
in sign, and a trimmed mean for the timers.

\section{Lessons from Earlier Harness Versions}
Briefly, why the design looks like this: position confounding, the mean dominated by a tail,
unpaired tiers, log writes in the measured path, the load generator competing with the edge.
\material{\repo{docs/BENCHMARK_METHODOLOGY.md} \S1}

%======================================================================
\chapter{RQ1: Latency Added}
\section{Median Latency per GET across Rates}
$+109$ to $+218$\us{} (13--23\%), resolved in every \texttt{GET} cell.
\material{\repo{results/paired_headline.md}, \repo{plots/paired_latency_added.png}}
\section{Tail Latency}
p95 grows more than the median on every edge. Apache at 100\,iterations/s collapses.
\section{Request Position and Reconnects}
The first request after the think-time sleep pays a wake-up cost. Apache closes the connection
after the trap's 500.
\material{\repo{docs/THESIS_NOTES.md} (2026-09-21, \S5b), \repo{plots/paired_latency_shift.png}}
\section{Answer to RQ1}
One paragraph: the number per edge, how it changes with load, and how certain it is.

%======================================================================
\chapter{RQ2: CPU and Capacity}
\section{CPU per Request}
CPU added per request: OpenResty $+33$, Envoy+Lua $+41$, WASM $+123$, Apache $+200$\us.
\section{Cost Attribution}
The in-edge timers cover 5--22\% of that CPU; the rest is VM entry and exit, body filters and log
writes.
\material{\repo{plots/cost_attribution.png}}
\section{The Probe Effect}
The instrumentation costs as much as WADM itself on two edges.
\section{Per-Mechanism Cost and Page Size}
The \texttt{GET} path against the trap; injection grows linearly up to 768\,KiB.
\material{\repo{results/mechanisms.md}, \repo{plots/size_sweep.png}}
\section{Capacity and the Binding Resource}
Share of bare capacity kept with silent WADM: OpenResty 71\%, WASM 71\%, Envoy+Lua 84\%, Apache
30\%, each reported with the resource that limited it.
\material{\repo{results/capacity_summary.md}}
\section{Answer to RQ2}
CPU and capacity per edge, and where the cost arises.

%======================================================================
\chapter{RQ3: Functional Correctness}
\section{Response Integrity under Load}
Share of intact responses per probe and rate. Apache is the only edge that falls below 100\%.
\material{\repo{results/paired_functional.md}}
\section{Detection Recall and False Positives}
\texttt{WADM ALERT} and \texttt{WADM TRAP} lines counted against expected; zero on bare.
\section{Response-Shape Matrix}
Gzip-encoded bodies and range requests are corrupted on all four edges. HEAD, bodies over 1\,MiB
and ETags each fail on some of them.
\material{\repo{results/functional_matrix.md}, \repo{benchmarks/functional_matrix.py}}
\section{Answer to RQ3}
What breaks, on which edge, and whether detection is complete.

%======================================================================
\chapter{RQ4: Influence of the Extension Runtime}
\section{Per-Kind Operation Cost}
In-edge timer distributions per honeytoken kind and edge.
\material{\repo{plots/paired_internal_timers.png}}
\section{The SQLi Trap: Scan Depth versus Decoding}
In the 2$\times$2, decoding costs 0.5--2.5\us{} on the Lua edges and nothing on WASM; scan depth
is visible only on PUC Lua.
\material{\repo{docs/THESIS_NOTES.md} (2026-09-21), \repo{results/sqli.md}}
\section{Case Study: Apache mod\_lua Segfaults}
Crashes driven by cumulative requests, and the \texttt{LuaScope conn} trade-off.
\material{\repo{docs/BENCHMARK_METHODOLOGY.md} \S4b--4c}
\section{Answer to RQ4}
How much of RQ1--RQ3 is explained by the runtime rather than by the deception logic.

%======================================================================
\chapter{Discussion}
\section{Synthesis across Research Questions}
One table of edges against the four answers.
\section{Recommendations for Operators}
Which edge to use, and which response types to exclude from injection.
\section{Threats to Validity}
\subsection{Construct Validity}
Whether silent $-$ bare really is ``the cost of WADM''; the timer regions.
\subsection{Internal Validity}
Host freezes, position effects, warm-up, load-generator contention.
\subsection{External Validity}
One VM, synthetic traffic, a static origin, no TLS.
\subsection{Conclusion Validity}
Five replicates, ranges rather than significance tests, resolution limits.

%======================================================================
\chapter{Conclusion}
\section{Summary of Answers}
RQ1--RQ4 in one sentence each.
\section{Future Work}
POST-body inspection on all four edges, compressed bodies, real attacker traffic, TLS and HTTP/2.

%======================================================================
\appendix
\chapter{Configuration Reference}
Every \texttt{config.json} field per honeytoken kind, and the SQLi trap's fields.
\material{\repo{CONTEXT.md}}
\chapter{Experiment Parameters and Run Log}
Suite IDs, rejected cells and their reasons, incidents such as the host sleeping mid-run.
\material{\repo{docs/BENCHMARK_METHODOLOGY.md} \S8}
\chapter{Full Result Tables}
Per-rate tables and figures.
\material{\repo{benchmarks/results/plots/by_rate/}}
\chapter{Reproduction Guide}
Commands and machine configuration, with how long each step takes.
\material{\repo{docs/BENCHMARK_METHODOLOGY.md} \S3--4}

\bibliographystyle{ieeetr}
\bibliography{sources}
\end{document}
